\section{September 18th: The GCD and The Euclidean Algorithm}

\subsection{The GCD}

\begin{defbox}
   Let $a$, $b \in \ZZ$ such that $a$ and $b$ are nonzero. The \textbf{greatest common divisor} (often abbreviated \textbf{GCD}) \textbf{of $a$ and $b$}, denoted $\gcd(a,b)$ is the largest positive integer $d$ such that $d \mid a$ and $d \mid b$. 
\end{defbox}

\begin{thmbox}
   Let $a$, $b \in \ZZ$ nonzero. Let $d := \gcd(a,b)$. Then every common divisor of $a$ and $b$ also divides $d$. Furthermore, there exists $u$, $v \in \ZZ$ such that:
   \[ d = ua + vb \]
\end{thmbox}
Moreover, $d$ is the smallest positive integer of the form $d = ua + vb$. 
\begin{proof}
   We first prove $d = ua + vb$. Let $S := $.
\end{proof}
